
\documentclass{article}
\usepackage{colortbl}
\usepackage{makecell}
\usepackage{multirow}
\usepackage{supertabular}

\begin{document}

\newcounter{utterance}

\centering \large Interaction Transcript for game `eqbench', experiment `eq\_bench', episode 62 with reflexion{-}llama{-}3.1{-}8b{-}instruct.
\vspace{24pt}

{ \footnotesize  \setcounter{utterance}{1}
\setlength{\tabcolsep}{0pt}
\begin{supertabular}{c@{$\;$}|p{.15\linewidth}@{}p{.15\linewidth}p{.15\linewidth}p{.15\linewidth}p{.15\linewidth}p{.15\linewidth}}
    \# & \multicolumn{2}{c}{Player} && \multicolumn{2}{c}{Game Master} \\
    \hline

    \theutterance \stepcounter{utterance}  
    & & & \multicolumn{4}{p{0.6\linewidth}}{
        \cellcolor[rgb]{0.9,0.9,0.9}{
            \makecell[{{p{\linewidth}}}]{
                \texttt{\tiny{[P1$\langle$GM]}}
                \texttt{You have attempted this same task before and did not succeed. These are the plans you wrote after those attempts.} \\
\texttt{Trial \#1: New plan: } \\
\\ 
\texttt{Upon re{-}reading the dialogue, I realized that I focused too much on the surface{-}level conversation and didn't consider the underlying dynamics between Marcus and Julia. I should have paid more attention to the power struggle and the tension between them. Julia's responses suggest that she's not intimidated by Marcus's warnings and threats, and she's willing to take risks to uncover the truth. This indicates that she's feeling confident and assertive, rather than grateful or relieved.} \\
\\ 
\texttt{To improve my analysis, I'll focus on identifying the key themes and emotions in the dialogue, such as the power struggle, the tension between truth and lies, and Julia's determination to uncover the truth. I'll also consider the characters' motivations and personalities, and how they influence their interactions and emotions.} \\
\\ 
\texttt{I'll give the following scores:} \\
\\ 
\texttt{Grateful: 0} \\
\texttt{Suspicion: 9} \\
\texttt{Relief: 0} \\
\texttt{Annoyance: 1} \\
\\ 
\\ 
\\ 
\texttt{Your task is to predict the likely emotional responses of a character in this dialogue:} \\
\\ 
\texttt{Marcus: Stuck in an elevator. Just my luck.} \\
\texttt{Julia: Could be worse. Could be stuck in a room with no windows.} \\
\texttt{Marcus: You're an optimist.} \\
\texttt{Julia: Realist. I just choose to see the glass half full.} \\
\texttt{Marcus: And what if the glass is poisoned?} \\
\texttt{Julia: Then it's a good thing I'm thirsty.} \\
\texttt{Marcus: You always this reckless?} \\
\texttt{Julia: Only when it matters.} \\
\texttt{Marcus: And what matters to you?} \\
\texttt{Julia: The truth.} \\
\texttt{Marcus: Truth is a dangerous game.} \\
\texttt{Julia: Only for those with something to hide.} \\
\texttt{Marcus: Or those who seek it.} \\
\texttt{Julia: I'm not afraid.} \\
\texttt{Marcus: You should be.} \\
\texttt{Julia: Is that a threat?} \\
\texttt{Marcus: A warning.} \\
\texttt{Julia: What do you know?} \\
\texttt{Marcus: More than you can handle.} \\
\texttt{Julia: Try me.} \\
\texttt{Marcus: I've seen good people destroyed by the truth.} \\
\texttt{Julia: And I've seen bad people empowered by lies.} \\
\texttt{Marcus: You're playing with fire, Ms. Winston.} \\
\texttt{Julia: Then I'll get burned.} \\
\texttt{Marcus: Not if I can help it.} \\
\texttt{{[}End dialogue{]}} \\
\\ 
\texttt{At the end of this dialogue, Julia would feel...} \\
\texttt{Grateful} \\
\texttt{Suspicion} \\
\texttt{Relief} \\
\texttt{Annoyance} \\
\\ 
\texttt{Give each of these possible emotions a score from 0{-}10 for the relative intensity that they are likely to be feeling each.} \\
\\ 
\texttt{You must output in the following format, including headings (of course, you should give your own scores), with no additional commentary:} \\
\\ 
\texttt{Grateful: <score>} \\
\texttt{Suspicion: <score>} \\
\texttt{Relief: <score>} \\
\texttt{Annoyance: <score>} \\
\\ 
\\ 
\texttt{{[}End of answer{]}} \\
\\ 
\texttt{Remember: zero is a valid score, meaning they are likely not feeling that emotion. You must score at least one emotion > 0.} \\
\\ 
\texttt{Your answer:} \\
            }
        }
    }
     \\ \\

    \theutterance \stepcounter{utterance}  
    & \multicolumn{4}{p{0.6\linewidth}}{
        \cellcolor[rgb]{0.561,0.737,0.561}{
            \makecell[{{p{\linewidth}}}]{
                \texttt{\tiny{[P1$\rangle$GM]}}
                \texttt{Grateful: 0} \\
\texttt{Suspicion: 9} \\
\texttt{Relief: 0} \\
\texttt{Annoyance: 6} \\
            }
        }
    }
    & & \\ \\

    \theutterance \stepcounter{utterance}  
    & & & \multicolumn{2}{p{0.3\linewidth}}{
        \cellcolor[rgb]{0.9,0.9,0.9}{
            \makecell[{{p{\linewidth}}}]{
                \texttt{\tiny{[GM$|$GM]}}
                \texttt{Grateful: 0} \\
\texttt{Suspicion: 9} \\
\texttt{Relief: 0} \\
\texttt{Annoyance: 6} \\
            }
        }
    }
    & & \\ \\

    \theutterance \stepcounter{utterance}  
    & & & \multicolumn{2}{p{0.3\linewidth}}{
        \cellcolor[rgb]{0.9,0.9,0.9}{
            \makecell[{{p{\linewidth}}}]{
                \texttt{\tiny{[GM$|$GM]}}
                \texttt{Grateful: 0} \\
\texttt{Suspicion: 8} \\
\texttt{Relief: 0} \\
\texttt{Annoyance: 4} \\
            }
        }
    }
    & & \\ \\

    \theutterance \stepcounter{utterance}  
    & & & \multicolumn{2}{p{0.3\linewidth}}{
        \cellcolor[rgb]{0.9,0.9,0.9}{
            \makecell[{{p{\linewidth}}}]{
                \texttt{\tiny{[GM$|$GM]}}
                \texttt{game\_result = LOSE} \\
            }
        }
    }
    & & \\ \\

    \theutterance \stepcounter{utterance}  
    & & \multicolumn{4}{p{0.6\linewidth}}{
        \cellcolor[rgb]{0.95,0.95,0.95}{
            \makecell[{{p{\linewidth}}}]{
                \texttt{\tiny{[P1$\leftrightarrow$P]}}
                \texttt{Post{-}episode Reflexion (not sent to GM):} \\
\texttt{New plan: } \\
\\ 
\texttt{Upon re{-}reading the dialogue, I realized that I focused too much on the surface{-}level conversation and didn't consider the underlying dynamics between Marcus and Julia. I should have paid more attention to the power struggle and the tension between them. Julia's responses suggest that she's not intimidated by Marcus's warnings and threats, and she's willing to take risks to uncover the truth. This indicates that she's feeling confident and assertive, rather than grateful or relieved.} \\
\\ 
\texttt{To improve my analysis, I'll focus on identifying the key themes and emotions in the dialogue, such as the power struggle, the tension between truth and lies, and Julia's determination to uncover the truth. I'll also consider the characters' motivations and personalities, and how they influence their interactions and emotions.} \\
\\ 
\texttt{I'll pay closer attention to the following:} \\
\\ 
\texttt{1. Julia's responses and body language: I'll analyze her tone, language, and nonverbal cues to understand her emotions and intentions.} \\
\texttt{2. Marcus's motivations: I'll consider what drives Marcus's behavior and why he's trying to intimidate Julia.} \\
\texttt{3. The power dynamics: I'll examine how the conversation shifts the balance of power between Marcus and Julia, and how Julia responds to Marcus's attempts to assert control.} \\
\texttt{4. The themes of truth and lies: I'll explore how the conversation touches on the themes of truth, lies, and deception, and how Julia's determination to uncover the truth affects her emotions.} \\
\\ 
\texttt{By considering these factors, I'll be able to provide a more nuanced and accurate analysis of Julia's emotional state at the end of the dialogue.} \\
            }
        }
    }
    & \\ \\

\end{supertabular}
}

\end{document}
